
\section{Results and Interpretation}
\label{section:results}

%\subsubsection*{Key Findings}

The key coefficient in \eqref{eq:01}, $ \beta $, captures the impact of the adoption of MTI on salary and wage employment. Table \ref{tab:mainreg} shows that MTI increases salary employment, but not wage employment. Excluding the Amhara sample does not change the outcome, suggesting that the results are driven by the changes that occurred in states with previously suppressed languages. 

The conclusion that salary employment increases with MTI is consistent with the institutional and policy changes the country has experienced in the last three decades. According to the World Bank, the public sector is the largest employer of non-farm labor in the country, followed by the service, manufacturing and construction sectors.\footnote{For more detailed information on the distribution of salary and wage employment in Ethiopia, see \url{https://openknowledge.worldbank.org/bitstream/handle/10986/32093/Ethiopia-Employment-and-Jobs-Study.pdf?sequence=1&isAllowed=y}} The public sector is the largest employer of salaried jobs, with the education and health sectors leading the way, and most public services are provided by state governments that do not require the mastery of ethnic languages for employment. 

 %most service, manufacturing and construction jobs are private-sector wage employments

However, most wage employments are provided by the private sector, where there are no language requirements in the recruitment and retention of workers. Amharic is still the key language of communication in many urban centers in the country, thus fluency in the language is important to secure wage employments. With students educated in their native tongues losing their proficiency in Amharic \citep{taye_medium_2019}, therefore, it is not surprising that they are at a disadvantage when it comes to obtaining wage employment.  

\subsection*{Likely Channels}

The introduction of MTI can improve employment outcomes if its positive effects on school performance outweigh its negative effects due to reduced proficiency in a historically dominant language. We can explore which channel is the most prominent in Ethiopia, testing whether changes in test scores due to MTI are associated with labor market outcomes, employing an instrumental variables strategy, assuming that an increase in MTI intake in primary schools leads to improved educational outcomes. In this set-up, the interaction of Predicted MTI and MTI intake is used as an instrument for test scores measuring cognitive ability. 

Consistent with emerging empirical evidence on the impact of MTI on educational outcomes, we document that cognitive ability improves with increased MTI intake (Table \ref{tab:rev}). In addition, prospects for salary employment increase with gains in cognition due to MTI, suggesting that the positive effects of educational gains due to MTI have exceeded the negative effects of the policy on the mastery of the dominant language in the market for salary employment. However, despite the reasonably well-established positive link between improved cognition and increased employment opportunities, elevated test scores due to increased MTI intake do not seem to improve wage employment outcomes, further supporting the conclusion that the treatment has reduced employability in the private sector by reducing proficiency in the dominant language. 

%The conclusion is in agreement with the reality that local governments that use the various ethnic languages as working languages are the largest employers of salaried jobs, although the mastery of Amharic is still essential for securing salaried jobs with the federal government.


In summary, the introduction of mother tongue education in Ethiopia has proven to be a double-edged sword for its beneficiaries. Although it increased their competitiveness on the market by raising their cognitive capacity, it has put them at a distinct disadvantage when it comes to securing private sector wage employment, which requires proficiency in Amharic, in which many have become less conversant due to the policy changes in languages of instruction. Potential workers in the private sector do not have legal rights to work in their native languages in Ethiopia, and there are no laws in the books that prohibit private employers from requiring workers to be fluent in Amharic as a criterion for employment, restricting the employment opportunities of individuals whose educational outcomes have increased due to MTI. 

  



